\chapter{Geometric Primitives and Swept Geometries}
\label{chap:geometry}

\section{Geometric Primitive Factory Architecture}
Rather than relying on static external OBJ or FBX model files, all geometry in \textit{Matsuri Nights} is generated procedurally via mathematical functions defined in \texttt{src/Primitives.h} and \texttt{src/Curves.h}. This procedural architecture ensures that vertex normals, texture coordinates, and tangent frames are mathematically exact and eliminates external asset loading overhead.

\section{Standard Geometric Primitives}
Table~\ref{tab:primitives_summary} lists the fundamental primitives provided by \texttt{Primitives.h}, detailing their vertex layout, triangle counts, and procedural construction.

\begin{table}[H]
\centering
\small
\begin{tabularx}{\textwidth}{p{2.5cm} p{2.2cm} p{2.2cm} p{7.5cm}}
\toprule
\textbf{Primitive} & \textbf{Vertex Count} & \textbf{Index Count} & \textbf{Construction Methodology \& Normal Layout} \\
\midrule
\textbf{Cube} \newline (\texttt{createCube}) &
24 vertices &
36 indices \newline (12 triangles) &
6 independent quad faces ($+Z, -Z, -X, +X, +Y, -Y$). Each face possesses 4 unique vertices with duplicate corner positions to ensure sharp, unshared face normals ($\pm X, \pm Y, \pm Z$). \\
\addlinespace
\textbf{Cylinder} \newline (\texttt{createCylinder}) &
102 vertices &
288 indices \newline (96 triangles) &
Constructed with $S = 24$ radial segments. Comprises $(S+1)\times 2 = 50$ side wall vertices with smooth radial normals $\mathbf{N} = [\cos\theta, 0, \sin\theta]^T$, 26 top disc cap vertices with $\mathbf{N} = [0, 1, 0]^T$, and 26 bottom cap vertices with $\mathbf{N} = [0, -1, 0]^T$. \\
\addlinespace
\textbf{Cone} \newline (\texttt{createCone}) &
52 vertices &
144 indices \newline (48 triangles) &
Side mantle evaluated over $S = 24$ segments with slanted outward normals $\mathbf{N} = \operatorname{normalize}([r\cos\theta, r/h, r\sin\theta]^T)$, sharing apex vertex, plus a closed circular base cap disc. \\
\addlinespace
\textbf{Sphere} \newline (\texttt{createSphere}) &
525 vertices &
2880 indices \newline (960 triangles) &
UV sphere evaluated across $R = 20$ latitude rings and $S = 24$ longitude sectors. Radial outward normals $\mathbf{N} = \mathbf{P} / r$, with polar caps degenerating gracefully into triangle fans. \\
\addlinespace
\textbf{Plane} \newline (\texttt{createPlane}) &
4 vertices &
6 indices \newline (2 triangles) &
Horizontal quad spanning $[-w/2, +w/2] \times [-d/2, +d/2]$ with upward normal $\mathbf{N} = [0, 1, 0]^T$ and normalized UV mapping $[0, 1]^2$. \\
\bottomrule
\end{tabularx}
\caption{Summary of Standard Geometric Primitives (\texttt{src/Primitives.h}).}
\label{tab:primitives_summary}
\end{table}

\section{Parametric Swept Tubes and Curved Architectural Beams}
\subsection{Generalized Swept Tubes via Bishop Frames}
The engine provides a generalized swept surface generator, \texttt{Curves::createSweptTube}, which takes arbitrary 3D parametric trajectory functions $\mathbf{P}(t)$, tangent evaluators $\mathbf{T}(t)$, and variable radius functions $R(t)$. Using the Bishop frame parallel transport algorithm described in Section~\ref{eq:bishop_frame}, a circle of radius $R(t_i)$ is swept across longitudinal slices:
\begin{equation}
\mathbf{V}_{i, j} = \mathbf{P}(t_i) + R(t_i) \cdot \left[ \cos(\phi_j) \mathbf{N}_i + \sin(\phi_j) \mathbf{B}_i \right]
\label{eq:swept_vertex}
\end{equation}
where $\phi_j = \frac{2\pi j}{S_{\text{rad}}}$ and surface normals correspond to $\mathbf{N}_{\text{surf}} = \cos(\phi_j) \mathbf{N}_i + \sin(\phi_j) \mathbf{B}_i$.

\subsection{Curved Architectural Beams (\textit{Kasagi} and \textit{Shimaki})}
The prominent upward-flaring top beam (\textit{Kasagi}) and sub-beam (\textit{Shimaki}) of the Torii gate cannot be represented via straight cylinders or boxes. The engine implements \texttt{Curves::createCurvedBeam}, which sweeps an asymmetric trapezoidal cross-sectional profile along a cubic B\'ezier spine curve:
\begin{itemize}
    \item \textbf{Cubic B\'ezier Spine:} Control points $\mathbf{P}_0 = [-3.8, 6.4, 0]$, $\mathbf{P}_1 = [-1.5, 6.15, 0]$, $\mathbf{P}_2 = [1.5, 6.15, 0]$, $\mathbf{P}_3 = [3.8, 6.4, 0]$, providing an authentic Japanese shrine curve with pronounced upward flare at the eaves (\textit{Sori}).
    \item \textbf{Cross-Section:} 4 corner vertices defining a trapezoid wider at the top than the bottom, with chamfered rain gutters along the roof ridge.
\end{itemize}

\subsection{Sacred Shimenawa Catenary Ropes}
The sacred straw rope spanning the Torii gate pillars is generated via \texttt{Curves::createCatenaryRope} (444 vertices, 864 triangles). Using Equations~\ref{eq:catenary_code} and~\ref{eq:catenary_tan}, the rope maintains a natural gravitational sag with 36 longitudinal slices and 12 radial segments.

\section{Organic Botanical and Human Anatomical Meshes}
To avoid sterile blocky geometry, \texttt{src/Curves.h} implements specialized procedural synthesizers for botanical and human characters:
\begin{enumerate}
    \item \textbf{Curved Cherry Blossom Petals (\texttt{createCurvedPetalMesh}):} Generates individual petals featuring gentle 3D concave cupping along the central midrib and a distinctive notched petal apex tip.
    \item \textbf{Cherry Blossom Clusters (\texttt{createSakuraBlossomLobe}):} Assembles 5 radially arranged curved petals around a central pistil and stamen core, forming realistic blossom clusters.
    \item \textbf{Pine Needle Fans (\texttt{createPineNeedleClusterMesh}):} Generates radial needle tufts for bonsai trees and planter boxes.
    \item \textbf{Articulated Human Head (\texttt{createHumanHeadMesh}):} A contoured cranial mesh with tapered jawline, occipital curvature, and neck socket.
    \item \textbf{Kimono Torso (\texttt{createHumanTorsoMesh}):} Models traditional Japanese festival attire with chest taper, overlapping kimono lapels (\textit{Eri}), and a raised waist sash (\textit{Obi}).
    \item \textbf{Articulated Limbs (\texttt{createArticulatedLimbMesh}):} Tapered limb cylinders capped with rounded hemispheres representing shoulder, elbow, hip, and knee pivot sockets.
    \item \textbf{Geta Sandal Foot (\texttt{createGetaFootMesh}):} Generates traditional elevated wooden clog sandals (\textit{Geta}) featuring an elevated base plate, two vertical ground teeth (\textit{Ha}), and an angled fabric toe strap (\textit{Hanao}).
\end{enumerate}

\section{OpenGL Mesh Buffer Management}
Each synthesized primitive is uploaded to GPU memory via the \texttt{Mesh} class (\texttt{src/Mesh.h}). Vertex attributes are interleaved into a single Vertex Buffer Object (VBO):
\begin{lstlisting}[language=C++, basicstyle=\ttfamily\scriptsize]
struct Vertex {
    glm::vec3 Position;  // Offset 0,  Size 12 bytes
    glm::vec3 Normal;    // Offset 12, Size 12 bytes
    glm::vec2 TexCoords; // Offset 24, Size 8 bytes
}; // Total stride: 32 bytes
\end{lstlisting}
The Vertex Array Object (VAO) binds the VBO and Element Buffer Object (EBO), configuring layout locations:
\begin{itemize}
    \item \textbf{Location 0 (\texttt{aPos}):} \texttt{glVertexAttribPointer(0, 3, GL\_FLOAT, GL\_FALSE, 32, (void*)0)}
    \item \textbf{Location 1 (\texttt{aNormal}):} \texttt{glVertexAttribPointer(1, 3, GL\_FLOAT, GL\_FALSE, 32, (void*)12)}
    \item \textbf{Location 2 (\texttt{aTexCoords}):} \texttt{glVertexAttribPointer(2, 2, GL\_FLOAT, GL\_FALSE, 32, (void*)24)}
\end{itemize}
Draw calls are dispatched via indexed rendering: \texttt{glDrawElements(GL\_TRIANGLES, indexCount, GL\_UNSIGNED\_INT, 0)}.
